% =============================================================================
% Chapter 13: Conclusion -- Part II
% =============================================================================
\mychapter{13. Conclusion}

This part validated the thesis's automated XCT defect-detection pipeline on real PODFAM metal AM
samples, independently of the NIST CoCr benchmark used in Part I. Full-volume comparison confirmed,
on new data, the literature's finding that Bernsen local thresholding outperforms Otsu and Yen for AM
porosity segmentation, and additionally surfaced and diagnosed a specific calibration failure mode in
Bernsen's automatic contrast-threshold estimation, a failure mode subsequently reproduced, at a
different severity, on a sixth sample processed after the original validation. A 2D U-Net trained on
Bernsen pseudo-labels, using a defect-content-ranked rather than random sample assignment, achieved a
best validation Dice of 0.600 and, on a fully held-out sample, predicted a total defect volume within
0.002 percentage points of the pseudo-label's own value, a result that argues for reporting
volume-level agreement alongside pixel-level overlap metrics in future evaluation of this pipeline.
The validation set's small size and the absence of independently verified ground truth are stated
explicitly as limits on how far these results can be generalised, consistent with the standard set in
Part I's own treatment of its NIST-phase pseudo-labels.

Taken together with Part I, this document now covers method development on a public benchmark (NIST
CoCr), full-volume classical-vs-deep-learning validation on real industrial process-monitoring data
(PODFAM, five core samples), and an out-of-distribution generalisation check on a sixth sample that
both supported the trained detector's usefulness and surfaced a reproducible, previously undocumented
failure mode in one of its classical comparison baselines.
